%% notes-tex/microbe-perturb-seq/sections/2-what-perturb-seq-measures.tex
%% SOURCE: hand-authored prose (jargon primer). Per-cell recovery figures are
%% cross-referenced to Sec. 4, which derives them from experiments/019-perturb-seq-costing/.

\section{Measurement Principles and Terminology\secstatus{todo}}
\label{sec:primer}

This section fixes the vocabulary. Everything after it assumes these terms.

\subsection{Molecular Identifiers\secstatus{todo}}
\label{sec:identifiers}

% \Cref, not \cref: cleveref lowercases the type name, so the sentence-initial
% form has to be the capitalising variant or the page opens with "table 1 is...".
\Cref{tab:glossary} is the reference: every term the document uses in a
technical sense, with a link to the section that treats it. It is worth a first
pass rather than only being consulted, because the naming is genuinely
treacherous --- four distinct things in this document are called a ``barcode'',
two different things are called an ``index'', and the difference between cells
per \emph{target gene} and cells per \emph{guide} is a factor of six that changes
every budget in \cref{sec:application}. The prose below then develops the four
identifiers properly.

% \medskip either side: the glossary is a longtable, not a float, so it sits
% exactly here and gets no float separation of its own.
\medskip
%% GENERATED FILE -- do not hand-edit.
%% SOURCE: experiments/019-perturb-seq-costing/scripts/render_tex_tables.py
\begingroup\footnotesize
\setlength{\LTleft}{0pt}\setlength{\LTright}{0pt}
\begingroup\captionsetup{type=table}
\captionof{table}[]{Controlled vocabulary. Terms are grouped by what they are \emph{about} rather than alphabetised, so related ones read together; the last column links to the section that treats each in depth. Generated from \file{experiments/019-perturb-seq-costing/scripts/glossary.py}.}\label{tab:glossary}
\endgroup
\begin{longtable}{@{}>{\raggedright\arraybackslash}p{0.21\textwidth}>{\raggedright\arraybackslash}p{0.62\textwidth}>{\raggedright\arraybackslash}p{0.11\textwidth}@{}}
\toprule
Term & Definition & See \\
\midrule
\endfirsthead
\multicolumn{3}{@{}l}{\footnotesize\itshape Table~\ref{tab:glossary}, continued}\\
\toprule
Term & Definition & See \\
\midrule
\endhead
\bottomrule
\endlastfoot
\multicolumn{1}{@{}>{\raggedright\arraybackslash}p{0.21\textwidth}@{}}{\bfseries Tags on a molecule} & & \\*[1pt]
\textbf{Cell barcode}\newline {\color{tcgray}(CB)} & A sequence shared by every cDNA molecule from one cell and ideally unique to it; what makes the data single-cell rather than bulk. Delivered pre-synthesised on a bead in droplet methods, built up across rounds of well-specific ligation in split-pool. & \cref{sec:identifiers} \\
\textbf{Unique molecular identifier}\newline {\color{tcgray}(UMI)} & A short random sequence attached to each cDNA molecule \emph{before} amplification. Reads sharing CB, gene and UMI are PCR copies of one original mRNA, so collapsing on the UMI turns read counts into molecule counts. 10\,nt in the split-pool protocols here, 10--12\,nt in droplet. & \cref{sec:identifiers} \\
\textbf{Sublibrary index} & The Illumina index read added at final PCR. In split-pool it is not merely a sample tag: cells are divided into sublibraries \emph{after} the last in-cell barcoding round, so the index is a further barcode dimension and two cells that took identical plate paths are still resolved if they landed in different sublibraries. This is why unique dual indices are mandatory -- an index hop manufactures a barcode collision. & \cref{sec:collisions} \\
\textbf{Guide barcode}\newline {\color{tcgray}(GBC)} & An expressed, polyadenylated sequence that stands in for the guide RNA itself. A Pol\,III guide cassette carries no poly(A) tail and is therefore invisible to 3$'$ capture, so the perturbation identity has to be re-encoded in something the assay can see. & \cref{sec:guide-capture} \\
\textbf{Barcode collision} & Two distinct cells receiving the same full barcode and being merged into one apparent cell. Rate is $1-((B-1)/B)^{C-1}$ for $C$ cells in a barcode space of size $B$, so it is controlled by adding rounds or sublibraries, not by sequencing harder. & \cref{sec:collisions} \\
\textbf{UMI collision} & Two distinct molecules of the same gene in the same cell drawing the same UMI, and being counted once. Negligible at yeast expression levels: a 10\,nt UMI has $4^{10}\approx10^{6}$ states against at most a few thousand molecules of any one gene. Not to be confused with barcode collision, which is a cell-level error. & \cref{sec:identifiers} \\
\addlinespace[3pt]
\multicolumn{1}{@{}>{\raggedright\arraybackslash}p{0.21\textwidth}@{}}{\bfseries Cell isolation} & & \\*[1pt]
\textbf{Droplet isolation} & One cell co-encapsulated with one barcoded bead in an emulsion droplet; the droplet wall is what keeps a cell's molecules together. Reagent cost scales linearly with cells, because you pay per channel. & \cref{sec:isolation} \\
\textbf{Combinatorial split-pool barcoding} & No isolation at all. Fixed, permeabilised cells act as their own reaction vessels; the population is split across a plate, barcoded in situ, pooled, and re-split. After $R$ rounds of 96 wells there are $96^{R}$ possible paths. Cost scales per well, not per cell, so a plate costs the same for ten thousand cells or a million. & \cref{sec:isolation} \\
\textbf{Doublet} & Two cells reported as one. In droplet methods it is two cells in one partition, governed by Poisson loading (10x quote $\sim$8\% at 20{,}000 cells recovered); in split-pool the analogous failure is a barcode collision, which has a different and controllable cause. & \cref{sec:isolation} \\
\textbf{Spheroplast} & A yeast cell whose $\beta$-glucan wall has been digested away, here with zymolyase under osmotic support. Required because no commercial chemistry can get reagents through an intact fungal wall; the trade is that a wall-less cell bursts in a normal buffer. & \cref{sec:why-hard} \\
\addlinespace[3pt]
\multicolumn{1}{@{}>{\raggedright\arraybackslash}p{0.21\textwidth}@{}}{\bfseries Split-pool chemistry} & & \\*[1pt]
\textbf{Linker strand} & A splint oligo pre-annealed to a barcode, positioning it against the 5$'$ end of the growing strand so T4 DNA ligase can seal the nick. Ligation is templated by the linker, which is why a stray barcode from an earlier well would corrupt the address. & \cref{sec:brettner} \\
\textbf{Blocking strand} & Added after each ligation round to outcompete the linker and retire that well's barcode, so it cannot participate in a later round. Without it, leftover oligos scramble the barcode sequence. & \cref{sec:brettner} \\
\textbf{Template switch} & A bead-borne reaction adding a common handle to the far end of the cDNA, giving the molecule PCR primer sites at both ends. Everything after it is conventional Illumina library prep. & \cref{sec:brettner} \\
\textbf{Source plate / working plate} & A source plate is the pre-diluted barcode-oligo stock as shipped; working plates are aliquots cycled in daily use. The distinction is economic: source plates are retired by freeze--thaw damage rather than by depletion, so they behave as a capital asset. & \cref{sec:startup} \\
\textbf{Protocol run} & One complete pass through the barcoding workflow -- fix, permeabilise, three rounds, split into sublibraries. Consumes one withdrawal from each barcode plate and carries up to $\sim$480{,}000 cells. A unit of \emph{bench work}, not of sequencing: the sublibraries it yields are sequenced and priced separately. & \cref{sec:startup} \\
\textbf{Sublibrary} & An aliquot of barcoded cells (5{,}000--20{,}000) split off after the last in-cell round, lysed and prepared as one Illumina library. Sublibrary count is set jointly by barcode-collision targets and by PCR complexity limits. & \cref{sec:startup} \\
\textbf{rRNA depletion} & Removal of ribosomal cDNA from the amplified library, after barcoding and before fragmentation and indexing. It does not improve capture; it stops $\sim$94\% of purchased reads being spent on rRNA, and is the single highest-leverage change available to a split-pool screen. & \cref{sec:rrna} \\
\addlinespace[3pt]
\multicolumn{1}{@{}>{\raggedright\arraybackslash}p{0.21\textwidth}@{}}{\bfseries Screen design} & & \\*[1pt]
\textbf{Perturb-seq} & A pooled genetic screen read out by single-cell RNA-seq, so each cell yields both a perturbation identity and a whole-transcriptome phenotype. The transcriptome, not a growth rate, is the phenotype. & \cref{sec:primer} \\
\textbf{Plex}\newline {\color{tcgray}($k$)} & Guide RNAs delivered per cell, so $k$ genes are perturbed in that cell at once. Raising $k$ divides the cell requirement for main effects by $k$; it does not make genome-wide pairwise interactions reachable. & \cref{sec:multiplex} \\
\textbf{Main effect} & The average effect of perturbing one gene, pooled over every cell carrying a guide against it regardless of what else that cell carries. Well defined only under an additive (log-additive) model. & \cref{sec:multiplex} \\
\textbf{Guide-pooling vs cell-pooling} & Two ways to compress a screen. Guide-pooling puts several guides in one cell; cell-pooling combines separately perturbed cells before sequencing. Both rely on the perturbations composing additively, and both are decoded rather than read directly. & \cref{sec:yao} \\
\textbf{FR-Perturb} & The decoder for a compressed screen: factorises the observed expression matrix by sparse PCA, then recovers per-perturbation effects by LASSO against the pooling design. The recovered effects are estimates, so power depends on the design matrix as much as on cell count. & \cref{sec:yao} \\
\textbf{Cells per target gene} & Cells carrying a guide against that gene, pooled over all of its guides. \emph{Not} cells per guide and \emph{not} cells per plasmid -- the distinction is a factor of six here and is the most common way a screen budget is misread. & \cref{sec:cells-per-pert} \\
\addlinespace[3pt]
\multicolumn{1}{@{}>{\raggedright\arraybackslash}p{0.21\textwidth}@{}}{\bfseries Counting and cost} & & \\*[1pt]
\textbf{Shot noise} & The irreducible randomness of counting. Sequencing samples molecules at random, so a gene at true rate $\mu$ is observed with variance $\mu$ and relative error $1/\sqrt{\mu}$. A property of the measurement, not of the cell -- and the only noise that buying more reads removes. Also called counting or sampling noise. & \cref{sec:limits} \\
\textbf{Biological overdispersion}\newline {\color{tcgray}($\varphi_j$)} & Variance in a gene's expression between genetically identical cells in excess of shot noise, from cell-cycle phase, metabolic state and transcriptional bursting. Reading each cell more deeply does not reduce it; only more cells do. Equal to the negative-binomial dispersion, so $\sqrt{\varphi_j}$ is the biological coefficient of variation reported in the RNA-seq literature. & \cref{sec:design-equation} \\
\textbf{Depth sufficiency}\newline {\color{tcgray}($d^{*}$)} & The per-cell depth $1/(\varphi_j p_j)$ at which shot noise and biological overdispersion contribute equally. Below it a screen is measurement-limited and deeper reads help; above it the requirement is within a factor of two of a floor no sequencing can cross. & \cref{sec:design-equation} \\
\textbf{Pseudobulk} & Summing UMIs across all cells sharing a perturbation, to estimate that perturbation's mean expression. It is what the whole design is powered on: single cells are too shallow individually, so the quantity that matters is cells $\times$ depth. & \cref{sec:cells-per-pert} \\
\textbf{Biological floor} & The minimum cells a perturbation needs before its mean is interpretable, independent of sequencing. Cell-cycle phase and metabolic state must be averaged out, and depth cannot substitute. The published field standard is 100 cells. & \cref{sec:cells-per-pert} \\
\textbf{Usable vs sequenced cell} & A usable cell survives QC \emph{and} carries an identified perturbation; a sequenced cell is one whose reads were paid for. The two differ by $\sim$4$\times$ in split-pool, which is why a quoted cost per cell understates the real one. & \cref{sec:cost-per-cell} \\
\textbf{Read pair} & One paired-end fragment, the unit sequencing is priced in here. Both method families need paired-end runs -- one read carries the cDNA and the other the barcodes -- but they put the barcode on opposite reads. & \cref{sec:sequencing} \\
\textbf{PhiX spike-in} & A fraction of the loaded library that is control phage DNA. At $\ge$10\% it is a split-pool protocol requirement rather than hygiene: read 2 carries fixed linker sequence at defined cycles, so base diversity collapses and patterned flow cells mis-register clusters without it. Budget it as a tax on usable reads. & \cref{sec:sequencing} \\
\end{longtable}
\endgroup

\medskip

A single-cell RNA-seq read is only useful if you can answer three questions about
it: which cell did it come from, was it a distinct original molecule or a PCR
copy, and which gene is it. Different chemistries answer these differently, but
all of them attach the same four classes of tag.

\notesec{Cell barcode (CB)} A sequence shared by every cDNA molecule from one
cell and, ideally, unique to that cell. It is what makes the data
\emph{single-cell} rather than bulk. In droplet methods the CB is delivered
pre-synthesised on a bead; in split-pool methods it is \emph{built up} across
several rounds of well-specific ligation, so the barcode records the trajectory
a cell took through a series of plates (Sec.~\ref{sec:brettner}).

\notesec{Unique molecular identifier (UMI)} A short random sequence --- 10--12 nt
in the methods here --- added to each cDNA molecule \emph{before} any
amplification. Two reads sharing a CB, a gene, and a UMI are copies of one
original mRNA; two reads sharing CB and gene but differing in UMI are two
molecules. Collapsing on the UMI converts read counts, which are distorted by
PCR efficiency, into molecule counts, which are the actual measurement. The
count of distinct UMIs recovered from a cell is the single most important
quality number in this document: it is the denominator of every statistical
argument in Sec.~\ref{sec:limits}. \term{UMI collision} --- two distinct
molecules of the same gene drawing the same UMI --- is negligible at yeast
expression levels, since a 10 nt UMI has $4^{10}\approx10^{6}$ states against at
most a few thousand molecules of any one gene.

\notesec{Sample or sublibrary index} An Illumina index read added during the
final PCR. In a conventional experiment it separates samples sharing a flow
cell. In split-pool it does more than that, and the difference matters: because
cells are divided into sublibraries \emph{after} the last in-cell barcoding
round, the index is effectively an additional barcode dimension, so two cells
that took identical paths through the plates are still resolved if they landed
in different sublibraries (Sec.~\ref{sec:collisions}).

\notesec{Perturbation identifier} What makes the experiment a screen rather than
a survey. This is the hard part in yeast and is treated separately in
Sec.~\ref{sec:guide-capture}; the short version is that the guide RNA itself is
usually invisible to the assay, so a proxy has to be engineered.

\subsection{Cell Isolation Strategies\secstatus{todo}}
\label{sec:isolation}

\notesec{Physical isolation (droplet)} A microfluidic device co-encapsulates one
cell with one barcoded bead in an emulsion droplet. Every molecule tagged inside
that droplet shares a barcode because the droplet wall keeps it that way. Cell
loading follows Poisson statistics, so a low occupancy is used to keep
\term{doublets} --- two cells in one droplet, reported as one --- acceptably
rare; 10x quote about 8\% at 20{,}000 cells recovered. The cost of physical
isolation is that you pay per channel of reagent, and it scales linearly with
cells.

\notesec{Combinatorial split-pool barcoding} No isolation at all. Cells are
fixed and permeabilised so they become their own reaction vessels: their
membranes retain the cDNA while enzymes diffuse in. The population is
distributed across a 96-well plate, a well-specific barcode is attached in situ,
the cells are pooled, redistributed, and barcoded again. After $R$ rounds a cell
carries an ordered barcode $\mathrm{BC}_1\Vert\dots\Vert\mathrm{BC}_R$, and with
$96$ wells per round there are $96^R$ possible paths. Cells sharing round~1 are
overwhelmingly likely to differ by round~3. The consequence for cost is the
important one: \emph{you pay per well, not per cell}, so a plate costs the same
whether it holds ten thousand cells or a million.

\notesec{The landscape between the two poles} Those are the extremes;
\cref{fig:methods-map} places the six approaches actually used in yeast between
them and pairs each with the library chemistry it can support. The organising fact
is that the isolation choice and the barcode choice are not independent. A method
that pools cells before sequencing has to write identity into the molecule, and
where it lands is fixed by which oligo carries it. A method that keeps one cell per
well until the very end needs no cell barcode at all, because the Illumina sample
index already identifies the well, and the well already identifies the cell.

\tcfigfit[0.96]{figures/scrnaseq-method-taxonomy.pdf}{%
  \textbf{A cell barcode is only needed where cells are pooled before sequencing,
  and most of the method landscape follows from that one fact.}
  The yeast scRNA-seq landscape, reorganised from the review
  of~\citep{nadal-ribellesRiseSinglecellTranscriptomics2024}.
  \textbf{(a)} What keeps one cell's molecules separate, ordered by throughput and
  each labelled with the container that does the work. Combinatorial indexing is
  set apart because it is not an isolation method at all: the fixed cell is its own
  vessel. The band beneath splits the six by whether a per-cell phenotype is
  recorded, which is what decides whether a method can be paired with a reporter or
  a sorting gate.
  \textbf{(b)} The question that actually separates the library types, which is
  which molecule carries the cell barcode. There are four answers: the oligo-dT
  gives $3'$ libraries, the template-switch oligo gives $5'$, a barcode loaded into
  Tn5 gives the plate route used by the only $5'$ method yet run in yeast, and in
  full-length libraries nothing in the molecule carries it.
  \textbf{(c)} Where the tags end up, with all three library types drawn on one
  column grid in one orientation so they can be compared: $3'$ and $5'$ differ only
  in which end of the transcript sits beside the barcode, while full-length carries
  no barcode anywhere and is resolved instead by the sample index, one per well ---
  the same accounting that makes split-pool's sublibrary index barcode round~4
  (Sec.~\ref{sec:collisions}).
  Redrawn and reorganised from the text and Table~1 of that review rather than
  reproduced from its figure, and with the orientation fixed $5'$ to $3'$
  throughout, where the original mirrors its $3'$ panel. Drawn at true Nature print
  size in \file{notes/assets/drawio/scrnaseq-method-taxonomy.drawio}.}{fig:methods-map}

\subsection{Microbial Constraints on Single-Cell Chemistry\secstatus{todo}}
\label{sec:why-hard}

Every commercial single-cell chemistry was developed on mammalian cells, and
three properties of a yeast cell break its assumptions.

\begin{itemize}[leftmargin=1.2em,itemsep=1pt,topsep=2pt]
  \item \textbf{A cell wall.} Reagents cannot diffuse in and the cell cannot lyse
    on the mammalian schedule. Both method families solve this with zymolyase, a
    $\beta$-glucanase, but at different points in the workflow --- and \emph{when}
    you digest turns out to matter more than whether (Sec.~\ref{sec:jariani}).
  \item \textbf{Roughly a fifth the mRNA.} A yeast cell holds on the order of
    $10^{4}$ mRNA molecules against $10^{5}$ for a mammalian cell
    (Sec.~\ref{sec:transcript-content}). Less input means fewer UMIs recovered,
    which is why per-cell yeast data always looks poor next to a mammalian
    benchmark and why pooling cells is the measurement strategy rather than a
    convenience.
  \item \textbf{An overwhelming rRNA background.} About 85\% of a yeast cell's
    RNA is ribosomal. Any capture chemistry that is not strictly poly(A)-selective
    spends most of its reads on rRNA, and as Sec.~\ref{sec:application} shows,
    that single fraction dominates the sequencing bill.
\end{itemize}

\subsection{The Count Matrix and Estimand\secstatus{todo}}
\label{sec:output-object}

What comes out is a count matrix $\mathbf{Y}\in\mathbb{Z}_{\ge0}^{N\times G}$ over
$N$ cells and $G$ genes, plus an assignment $\mathbf{Z}\in\{0,1\}^{N\times L}$ of
cells to the $L$ library elements. Because per-cell depth is low, the analysis
does not treat a single cell as an observation of a perturbation's effect.
Instead cells sharing a guide are pooled into a \term{pseudobulk} profile, and
the estimand is

\begin{equation}
  \theta_{\ell j} \;=\; \log_2\frac{\mu_{\ell j}}{\mu_{\mathrm{ctrl},\,j}},
  \qquad \boldsymbol{\Theta}\in\mathbb{R}^{L\times G},
  \label{eq:theta}
\end{equation}

the log fold change of gene $j$ under element $\ell$ against non-targeting
controls. Two things follow, and they shape every later section. Single-cell
resolution is used for \emph{assignment and quality control}, not for the effect
estimate --- so the relevant precision question is how many cells share a guide,
not how deep any one cell was read. And $\boldsymbol{\Theta}$ is dense: a
genome-scale screen in one condition produces of order $6000\times6000$
coefficients, a directed map from every perturbation to every transcriptional
consequence. That object, not a hit list, is what makes the screen worth running
for model training.
